% AUTHOR: Amlal El Mahrouss
% PURPOSE: On Proofs.

\documentclass[10pt]{article}

%% ---- Core Packages ----------------------------------------
\usepackage[T1]{fontenc}
\usepackage[utf8]{inputenc}
\usepackage{lmodern}
\usepackage{microtype}

%% ---- Page Geometry ----------------------------------------
\usepackage[
top=1in, bottom=1in,
left=0.75in, right=0.75in,
columnsep=0.25in
]{geometry}

%% ---- Mathematics ------------------------------------------
\usepackage{amsmath, amssymb, amsthm, amsfonts}
\usepackage{mathtools}
\usepackage{bm}
\usepackage{textcomp}
\usepackage{esint}

%% ---- Figures & Tables -------------------------------------
\usepackage{graphicx}
\usepackage{booktabs}
\usepackage{multirow}
\usepackage{array}
\usepackage{caption}
\usepackage{subcaption}
\usepackage{float}

%% ---- Code Listings ----------------------------------------
\usepackage{listings}
\usepackage{xcolor}
\usepackage{algorithm}
\usepackage{algpseudocode}

\lstdefinestyle{codestyle}{
	backgroundcolor=\color{gray!8},
	basicstyle=\ttfamily\footnotesize,
	breakatwhitespace=false,
	breaklines=true,
	captionpos=b,
	commentstyle=\color{green!50!black},
	keywordstyle=\color{blue}\bfseries,
	stringstyle=\color{orange!70!black},
	numberstyle=\tiny\color{gray},
	numbers=left,
	numbersep=5pt,
	showstringspaces=false,
	frame=single,
	rulecolor=\color{gray!40},
	tabsize=2
}
\lstset{style=codestyle}

%% ---- Hyperlinks -------------------------------------------
\usepackage[
colorlinks=true,
linkcolor=blue!70!black,
citecolor=green!50!black,
urlcolor=blue!60!black
]{hyperref}
\usepackage{url}

%% ---- Bibliography -----------------------------------------
\usepackage[numbers, sort&compress]{natbib}

%% ---- Miscellaneous ----------------------------------------
\usepackage{enumitem}
\usepackage{siunitx}
\usepackage{cleveref}

%% ---- Theorem Environments ---------------------------------
\theoremstyle{plain}
\newtheorem{theorem}{Theorem}[section]
\newtheorem{lemma}[theorem]{Lemma}
\newtheorem{corollary}[theorem]{Corollary}
\newtheorem{proposition}[theorem]{Proposition}

\theoremstyle{definition}
\newtheorem{definition}[theorem]{Definition}
\newtheorem{example}[theorem]{Example}

\theoremstyle{remark}
\newtheorem{remark}[theorem]{Remark}

%% ---- Custom Commands --------------------------------------
\newcommand{\R}{\mathbb{Z}}
\newcommand{\N}{\mathbb{N}}
\newcommand{\Z}{\mathbb{Z}}
\newcommand{\E}{\mathbb{E}}
\newcommand{\Prob}{\mathbb{P}}
\newcommand{\norm}[1]{\left\lVert #1 \right\rVert}
\newcommand{\abs}[1]{\left\lvert #1 \right\rvert}
\newcommand{\inner}[2]{\left\langle #1,\, #2 \right\rangle}
\newcommand{\ie}{\textit{i.e.}\xspace}
\newcommand{\eg}{\textit{e.g.}\xspace}
\newcommand{\etal}{\textit{et al.}\xspace}

%% ---- Caption Formatting -----------------------------------
\captionsetup{
	font=small,
	labelfont=bf,
	format=hang,
	justification=justified
}

\renewcommand{\qedsymbol}{\ensuremath{\blacksquare}}

\usepackage{lettrine}

%% ---- Header / Footer --------------------------------------
\usepackage{fancyhdr}
\pagestyle{fancy}
\fancyhf{}
\fancyhead[R]{\small\thepage}
\renewcommand{\headrulewidth}{0.4pt}

%% ============================================================
%%  FRONT MATTER
%% ============================================================

\title{%
	\vspace{-1.5em}%
	\large\textbf{Using a Proof Assistant for Non-Trivial Theorems.}%
}
\author{By A. E. Mahrouss}
\date{\today}

%% ============================================================
\begin{document}
	%% ============================================================
	
	\maketitle
	\tableofcontents
	\newpage
	
	%% ============================================================
	\section{Abstract:}
	%% ============================================================
	
	\lettrine[lines=2, lhang=0.1, loversize=0.2]{P}ER previous documents, our duty is now to write proof statements in Lean.\\
	Lean is a formidable proof assistant. We could use it to prove smaller theorems to then prove bigger theorems.\\Or, practice using it (Lean) to learn our craft.

	%% ============================================================
	\section{Synopsis:}
	%% ============================================================

	%% ============================================================
	\subsection{Theorems and Proofs:}
	%% ============================================================

	\lettrine[lines=1, lhang=0.1, loversize=0.2]{L}ET the following theorem for $x, y, z \in \mathbb{Z}$. The following theorem per FLT (Fermat's Last Theorem) admits an identity for the following:
	
	\begin{theorem}
		There exists no solutions for $p^{n}$ in $\mathbb{Z}.$ If and only If:
		\begin{align*}
			x^{n} + y^{n} + z^{n} &= p^{n}, \quad \forall n, c, x, y, z, p \neq 0, n \geq 2.
		\end{align*}
		For all variables of $n, x, y, z, p \in \mathbb{Z}.$ We could go by induction per Fermat Last Theorem that $x^{n} + y^{n} + z^{n}$ equals to:	
		\begin{align*}
			x^{n} + y^{n} + z^{n} &= p^{n}.
		\end{align*}
		For $\forall x^{n} + y^{n} + z^{n} \in \mathbb{Z}.$:
		\begin{align*}
			x^{n} + y^{n} + z^{n} &= a^{n} + b^{n}.
		\end{align*}
		If and only If for $\forall x^{n} + y^{n} + z^{n} \in \mathbb{Z}.$:
		\begin{align*}
			x^{n} + y^{n} + z^{n} \neq 0, \quad x^{n} + y^{n} + z^{n} &\geq 2, \quad \forall x^{n} + y^{n} + z^{n} \in \mathbb{Z}.
		\end{align*}
		Thus we can induce the following for $\forall x^{n} + y^{n} + z^{n} \in \mathbb{Z}.$:
		\begin{align*}
			x^{n} + y^{n} + z^{n} -  a^{n} &= b^{n}, \quad x^{n} + y^{n} + z^{n} - b^{n} = a^{n}.
		\end{align*}
		If and only If for $\forall x^{n} + y^{n} + z^{n} \in \mathbb{Z}.$:
		\begin{align*}
			x^{n} + y^{n} + z^{n} -  a^{n} - b^{n} &= 0, \quad x^{n} + y^{n} + z^{n} \neq 0.
		\end{align*}
	\end{theorem}
	\begin{remark}
		There is a possibility that $x^{n} + y^{n} + z^{n} \geq 2$, may be replaced under some circumstances with $x^{n} + y^{n} + z^{n} \geq \mathbb{P}_{0}.$ Where $\mathbb{P}$ is a sequence of pair numbers greater than one.
	\end{remark}
	\subsection{Notes on Theorem 1.1:}
	\begin{remark}
		One method for proving Theorem 1.1 is using proofs by direct and/or induction of the FLT by Wiles.
	\end{remark}
	\section{Remarks:}
	\lettrine[lines=1, lhang=0.1, loversize=0.2]{P}ROOF sources have been distributed with the repository, one exercise would be for the reader to implement such non-trivial proof in Lean.\\However usage of Lean may not be warranted for every proof writing.
	\section{Going Forward:}
	\lettrine[lines=1, lhang=0.1, loversize=0.2]{W}HILE a short note, the purpose is for the development of methodology on using proof assistants for non-trivial conjectures and theorems of Pure Mathematics.
	
\nocite{*}
	
\bibliographystyle{acm}
\bibliography{A_Course_of_Pure_Mathematics, wiles, avigad, weil1967basic}
	
	
%% ============================================================
\end{document}
%% ============================================================
